\section{Συμπεράσματα}
Η παρούσα διερεύνηση έδειξε ότι η μέθοδος \textit{Matrix Pencil} μπορεί να χρησιμοποιηθεί αποτελεσματικά για την ταυτοποίηση ηλεκτρομηχανικών ταλαντώσεων από αποκρίσεις \textit{ringdown}. Η ποιότητα της μαθηματικής ανακατασκευής των σημάτων και η συνέπεια των εκτιμήσεων ως προς τα βιβλιογραφικά modes αναφοράς επιβεβαιώνουν ότι η συνολική μεθοδολογία είναι κατάλληλη για την ανάλυση της δυναμικής συμπεριφοράς των εξεταζόμενων συστημάτων.

Για το \textit{Kundur Two-Area System}, τα αποτελέσματα ήταν πιο ξεκαθαρά και πιο εύκολα στην ερμηνεία. Ο κυρίαρχος \textit{inter-area} mode εντοπίστηκε με συνέπεια κοντά στα $0.54$~Hz, ενώ οι \textit{intra-area} ταλαντώσεις εμφανίστηκαν στην περιοχή περίπου $1.0$--$1.1$~Hz~\cite{estimation2021}. Παράλληλα, η ανακατασκευή των σημάτων παρουσίασε πολύ υψηλή πιστότητα, γεγονός που ενισχύει την αξιοπιστία των εξαγόμενων ιδιοτιμών, και οι μέθοδοι \textit{k-Means} και \textit{k-Medoids} οδήγησαν σε συγκλίνοντα συμπεράσματα.

Για το \textit{IEEE 39-Bus System}, η ανάλυση ήταν πιο σύνθετη, όπως ήταν αναμενόμενο λόγω της μεγαλύτερης κλίμακας του συστήματος. Παρότι οι εκτιμήσεις εμφανίζουν μεγαλύτερη διασπορά, η ανάλυση παραμένει συμβατή με τα modes αναφοράς της βιβλιογραφίας~\cite{sarkar2024drga} και επιτρέπει την ανάδειξη των βασικών περιοχών ταλαντωτικής συμπεριφοράς ανά περιοχή ελέγχου.

Τέλος, το στάδιο \textit{data screening} και η συμπληρωματική εφαρμογή τεχνικών συσταδοποίησης αποδείχθηκαν χρήσιμα για την οργάνωση και την ερμηνεία των αποτελεσμάτων. Συνολικά, η ροή εργασίας από το \textit{PowerFactory} στην \textit{Python} παρέχει μια συνεπή βάση για περαιτέρω αυτοματοποιημένη ανάλυση της δυναμικής ευστάθειας σε συστήματα ηλεκτρικής ενέργειας.

\FloatBarrier

